\chapter{Background}
\label{cap:background}

This chapter presents the foundations needed to put the project in context. First, the clinical data generated during the hospitalization of adult patients and the concept of FTR are described. Then PM and the concept of event log are introduced. After that, PPM is described, together with its prediction types and the usual way of evaluating it. The chapter ends with a description of the five ML algorithms compared in this work.

\section{Clinical Data of Inpatients and Failure To Rescue}
\label{sec:datos_clinicos}

During the hospitalization of an adult patient, the hospital information system records the episode from admission to discharge. Associated with the episode are the administrative data (admission and discharge dates, admission type, responsible service, discharge reason), the demographic data of the patient, the admission diagnosis and the diagnostic coding at discharge according to the International Classification of Diseases (ICD-10), and the care acts that take place during the stay. These acts include surgical interventions, with their type, anesthesia and postoperative destination; inter-service consultations; transfers between units, in particular to the ICU; nursing assessments; vital-sign readings; and laboratory results.

In the nursing field, the assessment of the inpatient relies on standardized indices that measure fall risk, functional capacity, pressure injury risk and comorbidity. The VALENF instrument \cite{lunaaleixos2022}, developed and validated in adult hospitalization units, integrates that assessment into a single meta-instrument. It is one of the sources of the nursing assessments recorded in the project database.

The term FTR was coined by Silber et al.\ \cite{silber1992} to describe the death of a patient after a complication that, had it been detected and treated in time, might not have been fatal. Unlike the raw complication rate, which depends mainly on the patient's risk, FTR reflects the hospital's ability to recognize deterioration and respond to it, and it has been proposed as an indicator of quality of care \cite{ghaferi2009}. However, FTR is not well characterized in the medical literature: there is no consensus on which complications to consider, on the time horizon or on the precursors that announce it. In hospital information systems, those observable precursors are unplanned admission to the ICU, urgent consultations, prolonged length of stay and, ultimately, death. Analyzing the situations that may lead to an FTR from the data recorded during hospitalization is therefore a necessary first step toward characterizing it.

The care process studied in this work is the elective surgical admission. It takes place as follows. The patient is admitted to the hospital on a scheduled date to undergo a planned surgery. Before the surgery, the surgical team assesses the postsurgical risk. If the risk is high, an ICU bed is reserved for the patient. The patient then goes through the operating theater. At the end of the surgery, the patient leaves the theater with a planned postoperative destination: the ward, the ICU, the post-anesthesia care unit or another unit. During the stay, the clinical team may request consultations to other specialists. If a complication arises, the patient may be transferred to the ICU or to another hospital. Finally, the patient is discharged. In this process, FTR manifests itself as the difference between what was planned and what happened: the patient whose planned destination was the ward and who ends up in the ICU, or whose stay is prolonged well beyond what was expected. That distinction between \emph{planning} and \emph{evolution} is central to the experimental design of Chapter~\ref{cap:implementacion}. The reference model of this process, designed for this work, is presented in Section~\ref{sec:modelo}.

\section{Process Mining}
\label{sec:pm}

PM is the discipline that connects data science with Business Process Management (BPM). From the records left by information systems, it extracts knowledge about how processes are actually executed \cite{aalst2016}. Its starting point is the event log, a collection of events each of which belongs to a case (an instance of the process), corresponds to an activity and carries a timestamp. Events may also carry additional attributes, such as the resource that executes them, the service or the result of a decision. The ordered sequence of the events of a case is called a trace, and the set of distinct traces that appear in the event log, variants.

Three types of analysis are distinguished on an event log \cite{aalst2016}.

\begin{itemize}[leftmargin=*]
  \item \textbf{Process discovery}. It consists of building a model of the process from the event log, without prior information. The most widely used algorithms are the Heuristics Miner, which builds a causal net from the frequency dependencies between activities and tolerates noise, and the Inductive Miner, which discovers process trees guaranteeing well-formed models \cite{leemans2014}. The DFG is the simplest representation, with one node per activity and one arc per pair of consecutive activities, annotated with its frequency.
  \item \textbf{Conformance checking}. It consists of comparing the event log with a reference model, either normative or discovered, in order to measure how far the observed behavior matches the expected one. The reference technique is \emph{alignments}, which search, for each trace, the most similar path in the model and quantify the deviations \cite{carmona2018}. From them are derived \emph{fitness}, which measures what fraction of the behavior in the event log the model explains, and \emph{precision}, which measures what fraction of the behavior of the model appears in the event log.
  \item \textbf{Enhancement}. It consists of enriching the model with performance, resource or decision information.
\end{itemize}

In healthcare, PM has been applied to emergency flows, oncology processes and patient trajectories, with the particularity that clinical processes are highly variable, poorly structured and strongly patient-dependent \cite{munoz2022}. That variability is the reason why the choice of the scope and of the granularity of the event log is a central decision of the work.

\section{Predictive Process Monitoring}
\label{sec:ppm}

PPM is the branch of PM that, from the historical event log, builds models capable of predicting how a case that is still running will continue \cite{difrancescomarino2022}. The unit of prediction is the prefix, formed by the first $k$ events of a trace. For each prefix, something about the future of the case is to be predicted. Di Francescomarino et al.\ \cite{difrancescomarino2018} classify PPM methods into three types according to what is predicted:

\begin{itemize}[leftmargin=*]
  \item \textbf{Outcome-based predictions}: predictions related to predefined categorical or boolean outcome values of the case, for example whether it will end in a deadline violation, in a complaint or, in this work, in an ICU admission \cite{teinemaa2019}.
  \item \textbf{Numeric value predictions}: predictions related to measures of interest taking numeric or continuous values, such as the remaining time until the end of the case \cite{aalst2011} or its cost.
  \item \textbf{Next event predictions}: predictions related to sequences of future activities and their data payloads, of which the prediction of the next activity is the simplest case \cite{tax2017}.
\end{itemize}

This classification is used in the rest of the report to name the three predictions carried out. Prediction~1 is a next event prediction (the next activity of the case). Predictions~2 and~3 are outcome-based predictions with a boolean outcome: whether the patient will be admitted to the ICU and whether the length of stay will be prolonged. No numeric value prediction is addressed.

PPM approaches are traditionally divided into two families \cite{difrancescomarino2018}. Model-based approaches rely on an explicit representation of the process, such as a transition system or a stochastic Petri net, annotated with statistics from the event log. ML-based approaches encode prefixes as feature vectors or as sequences and train a classifier or a neural network. Within this second family there are two ways of encoding a prefix \cite{teinemaa2019}. In aggregate encoding the prefix is summarized in a fixed-length vector with the count of each activity, the last activity, the elapsed times and the case attributes; it is the natural input of decision trees, ensembles and linear models. In sequence encoding the prefix is kept as a sequence of events, each represented by an embedding of its activity and by numerical features, padded up to a maximum length; it is the input of recurrent networks and Transformers.

The evaluation of a PPM model generally follows the protocol set by the benchmark of Teinemaa et al.\ \cite{teinemaa2019} for outcome prediction. Cases are split between training and test with a temporal split, so that the oldest cases train and the most recent ones evaluate. All prefixes of length $1 \ldots k_{\max}$ of each case are generated. Performance is measured with the area under the receiver operating characteristic (ROC) curve (AUC-ROC), which quantifies the ability of the model to rank positive cases above negative ones regardless of the decision threshold. The AUC-ROC is also reported per prefix length, because it is of interest to know how much the prediction improves as the case progresses. When the positive class is a minority, as happens with adverse clinical outcomes, the AUC-ROC can be high even if the model gets few positives right, because negatives dominate. For this reason it is complemented with the area under the precision--recall curve (AUC-PR), whose random baseline is the prevalence of the positive class \cite{saito2015}, and with the Brier score, which measures the calibration of the probabilities. For next-activity prediction, being a multiclass problem, accuracy and macro F1-score are used, against a baseline that predicts the most frequent transition.

\section{Machine Learning Algorithms Evaluated}

Five algorithms are compared, covering the two ML families usual in PPM: three tree-based models, which work on the aggregate encoding, and two sequential neural networks.

\paragraph{Decision Tree (DT).} A DT \cite{quinlan1986,breiman1984} recursively partitions the feature space by choosing at each node the variable and the threshold that best separate the classes. It is the most interpretable model, since the prediction is explained as a chain of conditions, but it tends to overfit and is unstable under small changes in the data. It is included as the interpretable reference model, as in Zare et al.\ \cite{zare2025}.

\paragraph{Random Forest (RF).} RF \cite{breiman2001} trains many trees on bootstrap samples of the training set, choosing at each node a random subset of variables, and averages their predictions. Averaging reduces the variance of the individual tree and yields well-calibrated probabilities, and it also provides a measure of variable importance. It is the reference ensemble in the benchmark of Teinemaa et al.

\paragraph{Extreme Gradient Boosting (XGBoost).} XGBoost \cite{chen2016} is an implementation of gradient boosting \cite{friedman2001} in which trees are built sequentially, each correcting the errors of the previous ensemble, with regularization and early stopping. It is usually the most accurate tabular model when its hyperparameters are tuned, but it is more sensitive to the configuration than RF.

\paragraph{Long Short-Term Memory (LSTM).} LSTM networks \cite{hochreiter1997} are recurrent networks with memory cells and gates that control what is remembered and what is forgotten, which allows them to capture dependencies along a sequence. Tax et al.\ \cite{tax2017} introduced them in PPM to predict the next activity and the remaining time, and since then they have been the reference Deep Learning (DL) model in the field.

\paragraph{Transformer.} The Transformer \cite{vaswani2017} replaces recurrence with a multi-head \emph{self-attention} mechanism that directly relates all the elements of the sequence to each other, with a positional encoding that preserves order. In PPM it has been applied both in its standard form (ProcessTransformer \cite{bukhsh2021}) and in variants that attend to each attribute separately, such as the Dynamic Attribute-Aware Transformer (DAW-Transformer) of Zare et al.\ \cite{zare2025}. In this work a standard Transformer encoder of reduced size is used, in accordance with the size of the event log.
